\section{道路上的国家}
\label{sec:early-tang-acceptance-scenes-four}

\subsection{虎牢以后的河北不是一纸降书}
621年虎牢之战后，\eventref{ST-621-HULAO}{唐取洛阳}改变了中原最重要的城市、仓储与渡口的归属。它没有替唐朝取得河北各地的兵源、乡里与地方首领。王世充降唐是可核的政治结果，刘黑闼此后在河北的再起同样在提醒朝廷：一座都城易手与地方社会接受新的征收、任命和军事秩序是两件事。\eventanchor{ST-622-LIU-HEITA}{ST-622-LIU-HEITA}\eventref{ST-622-LIU-HEITA}{刘黑闼再起（622）}的失败使唐重新控制主要州县，却不能使我们精确知道每一村落在战争、报复和重新登记中承担了什么。

唐廷面对的选择不是单纯的宽或严。过度依赖旧有地方力量，可能保留未来动员的风险；过快更换人事和征发资源，又可能破坏刚刚恢复的交通与生产。任命、赦免、赏赐和驻军都试图把军事胜利变成日常治理。史书偏重将领与中央的决定，墓志和地方材料零星照见个人流动，二者之间仍有大量不可见的基层协商。连续叙事应保留这段间隔：建国不是在618年完成的名号，而是多年在城门、仓房与县署中反复确认的过程。

\subsection{梁师都与北方边缘的供给问题}
628年梁师都覆亡，\eventanchor{ST-628-LIANG-SHIDU}{梁师都覆亡（628）}，唐朝才在北方较完整地结束隋末割据的一段余波。朔方与关中北缘的意义不只在地图上补齐一块区域。它关联着草原方向的军路、河套的水草、城镇的粮食与突厥关系；在太原、长安与边地之间，任何一个节点都可能影响军队能否行动。梁师都能够维持，与外部支持、地方资源和道路可用性有关；其覆亡也应从这些条件理解，而不是只视为唐军最终必然到达的终点。

边缘地区的居民不会只以“叛”“附”两种身份存在。有人为城防、运输或市场工作，有人因亲属和生计保持旧联系，有人随局势改变迁居。史料往往在他们影响朝廷安全时才留下记录，因而不能替他们设定统一动机。对唐朝而言，收复朔方降低了北方压力，却也需要重新安排守备、征收和贸易。与东突厥的交涉正在同一时期展开，这两条线互相影响而不彼此替代：内地的统一让朝廷腾出资源处理草原，草原的不确定又继续决定边地必须保持多少防御。

\subsection{灾荒、赈济与政治语言的限度}
灾异在唐代史书中常紧接着诏令、减免或君主自责出现。这样的排列能显示朝廷如何把天气、粮食和统治责任联系起来，不能直接证明某次灾害覆盖的范围，或每项赈济都送到了需要的人手中。旱、涝、蝗、疫与道路中断会以不同方式触及不同区域；一处仓储、河道或市场的条件，可能让相邻地区产生不同后果。将灾害写成某位君主德行的简单证明，会重复史书的修辞而错过物质世界的复杂性。

赈济也并非只有中央一道命令。仓粮如何调拨，谁负责发放，流民能否抵达城市或亲属网络，地方富户与寺院是否提供援助，都会改变实际结果。现有材料很少给出完整账本，尤其难以显示没有留下文字的家庭如何度过危机。谨慎的写法应区分可确认的诏令、减免和运输，与对社会后果的推测。即使不能精确统计，也能看清财政和行政为何必须在平时维持仓、路与记录：危机来临时，国家能力首先表现为能否把资源送到错误最少、最不可能通行的地方。

\subsection{松州以东的马匹和以西的消息}
边防的消息从来不只沿一条方向传播。松州驻军需要知道高原道路、部众活动与季节变化，长安则依赖边将、使者与驿传把局部判断转化为可供决策的文字。传递本身会筛选信息：将领可能强调威胁以请求粮饷，朝廷可能以礼仪词汇重述军事摩擦，后来的史家又会按战争或婚盟的结局组织材料。\eventref{ST-638-SONGZHOU-TIBET}{松州冲突}因此不是纯粹的战场事件，也是不同层级如何彼此理解消息的问题。

马匹、草料和道路的关系使这一点更具物质性。骑兵的行动要看牧地、季节和补给，山道的通过要看天气、向导和城寨，远行的使节也需要同样的停宿与安全。我们不能由某次战报推算所有军马，不能由一段道路遗存断言某年行军路线的每一站；但这些限制足以说明，外交、战争和婚盟并非三个彼此隔离的领域。它们都围绕谁能通过、谁能供给、谁能把消息带回而展开。

\subsection{高宗朝的官员怎样看见地方}
高宗时期的任命把许多官员从都城送往州县、边镇和新设地区。官衔记录能够标明制度位置，不能说明他们实际拥有多少人手或地方信任。一个刺史要依靠属官、书吏、仓吏、军人和地方中介，才可能处理讼狱、户籍、交通与税收；一个都护则要在更远距离上依赖译者、城镇首领和驻军。中央的制度语言提供共同的名称，地方的工作却总在不共同的地形、市场和社会关系中完成。

墓志中频繁出现的迁任，显示仕宦家庭如何跨越区域建立联系，也提醒我们其可见性远高于普通家庭。某人从关中到河西、从洛阳到岭南的履历，可以帮助定位道路和官僚流动；志文的嘉奖却不能直接当作地方治理成功的证据。将墓志同诏令、地理志和文书并读，能够提出更节制的问题：何种人才被派往何处，哪些地方需要反复更换官员，流动又怎样影响家族与地方社会？

\subsection{白江口后的海上讯息}
663年白江口战后，\eventref{ST-663-BAEKGANG}{白江口之战}的结果需要经过海峡两岸不同的政治记忆才会成为历史。唐朝记录关注联军胜利和半岛秩序，日本材料关注援助百济的失败及其后果，新罗叙事则有自身的统一与竞争框架。三者并读不是为了选出唯一的“正确国家视角”，而是为了看见同一片水域如何连接不同的军政目标。船只从港口出发，需要木材、人力、风向、粮食与导航；战争结束后，逃亡者、使者和消息仍会沿海路移动。

海上关系也不能用陆上朝贡词汇一概处理。港口既是礼物、人员和文本进入的门，也是征税、检查、疾病与军事警戒的地方。考古遗存能显示某些物资和航路的联系，不能让我们直接知道每件货物的主人或目的；编年史能保留使节与战事的日期，不能替所有水手和沿岸居民发言。正是这种材料层次，要求初唐的外部史同时写海峡的物质条件与各方政治叙事。
